\subsection{Task-conditioned structural quotients: erasure must be licensed, not merely useful}
\label{sec:tcsq}

A directional witness answers whether a structural object already extracted from one problem can be transported into another. A preceding problem is more basic: which distinctions in the source problem should enter that structural object at all? A raw representation may contain entity names, narrative order, coordinate choices, formatting conventions and other surface variables together with the relations that actually determine a registered quantity of interest. Treating all of these coordinates as equally load-bearing can make retrieval and transfer depend on irrelevant detail. Removing them indiscriminately is equally unsafe. We therefore introduce a \emph{task-conditioned structural quotient} as a derived, evidence-bound computational view rather than as a replacement for the original problem.

Let $P$ denote an immutable source problem representation, $q$ a frozen quantity of interest and $c$ a frozen context. A quotient proposal is a map
\[
Q_{q,c}:P\longrightarrow Z,
\]
where $Z$ retains a registered set of structural coordinates and identifies a complementary set of candidate nuisance coordinates to erase or conditionally erase. The map is intentionally target dependent: a coordinate can be irrelevant for one QoI and indispensable for another.

\begin{definition}[Task-conditioned structural quotient contract]
A TCSQ proposal for $(P,q,c)$ is a tuple
\[
\mathcal Q=(I^+,E^-,E^?,\Gamma,I_{\rm inv},B,O,R,F,L),
\]
where $I^+$ are preserved coordinates, $E^-$ are erased coordinates, $E^?$ are conditionally erased coordinates, $\Gamma$ contains equivalence generators or declared nuisance transformations, $I_{\rm inv}$ are invariants that must survive the reduction, $B$ are protected coordinates, $O$ are sufficiency obligations, $R$ records reconstruction bindings back to the source problem, $F$ are falsifiers, and $L$ are forbidden losses. The three coordinate sets $I^+,E^-,E^?$ must form a disjoint cover of the registered source coordinates, and every protected coordinate must lie in $I^+$.
\end{definition}

The proposal has no authority to certify its own sufficiency. A separately bound validation report assigns a verdict
\[
V_Q\in\{\mathrm{VALID\_EXACT},\mathrm{VALID\_APPROXIMATE},\mathrm{REJECT},\mathrm{CANNOT\_CHECK}\}.
\]
Materialization for solver use is permitted only when every registered sufficiency obligation has been verified, no obligation has failed or remains unknown, protected-coordinate checks pass, evidence is content-bound to the exact proposal and source hashes, and at least one declared validation route has run. An approximate quotient additionally requires an explicit error metric and tolerance. Thus failure to establish irrelevance is not converted into permission to erase.

\paragraph{Solution sufficiency.}
For an exact quotient, the strongest ideal is that the target solution functional factors through the quotient on the declared scope:
\[
F_{q,c}(P)=H_{q,c}\bigl(Q_{q,c}(P)\bigr).
\]
The implementation does not infer this equality from compression ratio, retrieval performance or model confidence. Instead it records finite, typed obligations whose verification route can be domain specific: exhaustive oracle comparison on a known world, metamorphic tests over nuisance orbits, formal proof, counterexample search, or an externally governed scientific evaluation. For approximate quotients, the factorization is replaced by a registered bound such as
\[
d\!\left(F_{q,c}(P),H_{q,c}(Q_{q,c}(P))\right)\le \epsilon,
\]
with the metric $d$ and tolerance $\epsilon$ frozen before use.

\paragraph{Reconstruction and original-problem verification.}
A quotient-side result is never promoted merely because it is correct on $Z$. A reconstruction operator
\[
R(s_Z;P,\mathcal Q)\longrightarrow s_P
\]
restores the bindings required by the original problem, after which the ordinary original-problem verifier must accept $s_P$. This separates three claims that generic compression can easily blur:
\[
\text{quotient valid}\neq\text{quotient solution valid}\neq\text{original problem solved}.
\]
A reconstruction failure or failed original-problem verification therefore closes the route without rewriting the source problem.

\paragraph{Derived view rather than destructive simplification.}
The canonical source remains immutable. The current implementation materializes an accepted quotient as a derived memory view whose erasure ledger is explicit and whose authority certificates are copied from a pinned canonical source rather than caller-supplied. The same validated view may compile a derived problem fibre for retrieval; absence of a coordinate from that fibre means only that the coordinate is computationally suppressed for $(q,c)$ under the registered certificate. It does not mean that the coordinate is false, scientifically absent, or irrelevant to another target.

This distinction allows TCSQ to compose with the directional transfer machinery. The intended pipeline is
\[
P\xrightarrow{\;Q_{q,c}\;}Z
\longrightarrow S_{q,c}
\xrightarrow{\;w_{A\to B}\;}
\text{candidate transfer},
\]
where the quotient decides which source distinctions are admitted into the task-conditioned structural view and the later witness decides whether those retained relations may be transported to a target. Protected coordinates, preserved invariants and forbidden losses are carried into the obstruction/transport layer; erased coordinates are never silently reconstructed as transfer evidence.

\subsubsection{Nearest work and the residual claim}
Task-relevant compression itself is not new. The information-bottleneck principle formalizes compression subject to preservation of information about a relevant variable \cite{tishby2000ib}; state-abstraction work studies representations that preserve action or value-relevant behavior \cite{huang2022actionsufficient,abel2020value}; recent decision-sufficient representation work directly optimizes representations for downstream decisions \cite{ye2026decision}; and prompt-compression systems such as LongLLMLingua and LLMLingua-2 remove context tokens while attempting to preserve downstream model performance \cite{jiang2024longllmlingua,pan2024llmlingua2}. Formal verification likewise has a long tradition of abstraction with counterexample-guided refinement rather than trusting a coarse abstraction unconditionally \cite{clarke2000cegar}.

More recent parents narrow the residual further. Context Codec explicitly represents source-grounded semantic atoms, protects critical commitments, measures compression soundness/completeness, studies round-trip recoverability and recommends conservative fallback when low-confidence or safety-critical information would be lost \cite{trukhina2026contextcodec}. Consequently, ``auditable compression with protected information and verification'' is not a TCSQ novelty claim. SLICEFORMER learns static program slices with dataflow-aware representations and constrained decoding \cite{he2026sliceformer}, while STACK performs state-aware compression of reasoning traces to reduce redundant reasoning without discarding task-relevant knowledge \cite{sui2026stack}. These works also rule out novelty claims based merely on removing irrelevant program/context/reasoning state.

The candidate residual here is therefore narrower and problem-solving specific. TCSQ asks whether a scientific problem object can be quotiented under a frozen QoI/context by an explicit structural equivalence/erasure contract, then used as the input to the same problem-fibre and directional-transfer machinery while preserving source lineage, protected coordinates, invariants and forbidden losses. Solver-side sufficiency is separate from compression fidelity, and a quotient-side solution must be reconstructed and verified against the exact original problem before any scientific promotion. The quotient is also subject to the framework's authority non-escalation rule: computational omission cannot mint, revoke or upgrade scientific authority. In the stronger empirical programme, Context-Codec-style verifiable compression, generic prompt compression, state abstraction, structure-without-erasure and mechanism-alignment controls are therefore comparators rather than strawmen.

Whether these additional obligations improve total task cost, nuisance robustness or cross-domain retrieval remains an empirical question. The present implementation plus known-world SQ-1/SQ-2 studies establish a fail-closed contract and controlled sufficiency/dependency-recovery results; they do not establish that a language model can discover valid quotients, that TCSQ is minimal, or that it outperforms strong compression/abstraction baselines on natural research tasks.
